% Partie : intuitionistic-logic
% Chapitre : introduction
% Section : syntax

\documentclass[../../../include/open-logic-section]{subfiles}

\begin{document}

\olfileid[fr]{int}{int}{syn}

\olsection{Syntaxe de la logique intuitionniste}

La syntaxe de la logique intuitionniste est la même que celle de la
logique propositionnelle. En logique propositionnelle classique, on
peut définir certains connecteurs au moyen des autres : par exemple,
$!A \lif !B$ par $\lnot !A \lor !B$, ou $!A \lor !B$ par
$\lnot(\lnot !A \land \lnot !B)$. Aussi les présentations de la logique
classique introduisent-elles souvent certains connecteurs comme des
abréviations de ces définitions. Ces définitions ne conviennent pas
en logique intuitionniste, avec deux exceptions : $\lnot !A$ peut
être défini, et l'est souvent, comme une abréviation de
$!A \lif \lfalse$. Dans ce cas, bien entendu, $\lfalse$ doit lui-même
rester primitif !{} De plus, on peut définir $!A \liff !B$, comme
en logique classique, par $(!A \lif !B) \land (!B \lif !A)$.

Les !!{formula}s de la logique propositionnelle intuitionniste
se construisent à partir de \emph{!!{propositional variable}s} et
de la constante propositionnelle~$\lfalse$, à l'aide de
\emph{connecteurs logiques}. Nous disposons des symboles suivants :

\begin{enumerate}
\item Un ensemble !!{denumerable}~$\PVar$ de !!{propositional variable}s
  $\Obj p_0$, $\Obj p_1$, \dots.
\item La constante propositionnelle de !!{falsity}~$\lfalse$.
\item Les connecteurs logiques : $\land$ (conjonction),
  $\lor$ (disjonction) et $\lif$ (implication).
\item Les signes de ponctuation : (, ), ainsi que la virgule.
\end{enumerate}

\begin{defn}[Formule]
\ollabel{defn:formulas} L'ensemble~$\Frm[L_0]$ des
\emph{!!{formula}s} de la logique propositionnelle intuitionniste
est défini inductivement comme suit :
\begin{enumerate}
\item $\lfalse$ est une !!{formula} atomique.

\item Toute !!{propositional variable}~$\Obj p_i$ est une
  !!{formula} atomique.

\item Si $!A$ et $!B$ sont des !!{formula}s, alors $(!A \land !B)$
  est !!a{formula}.

\item Si $!A$ et $!B$ sont des !!{formula}s, alors $(!A \lor !B)$
  est !!a{formula}.

\item Si $!A$ et $!B$ sont des !!{formula}s, alors $(!A \lif !B)$
  est !!a{formula}.

\tagitem{limitClause}{Rien d'autre n'est !!a{formula}.}{}
\end{enumerate}
\end{defn}

Outre les connecteurs primitifs introduits ci-dessus, nous utilisons
les symboles \emph{définis} suivants : $\lnot$ (négation) et $\liff$
(!!{biconditional}). Les formules contenant ces opérateurs définis
s'interprètent comme suit :
\begin{enumerate}
\item $\lnot !A$ abrège $!A \lif \lfalse$.

\item $!A \liff !B$ abrège $(!A \lif !B) \land (!B \lif !A)$.
\end{enumerate}

Bien que $\lnot$ soit officiellement une abréviation, nous donnerons
parfois des règles explicites ou des clauses de définition pour~$\lnot$,
comme s'il était primitif. Cela permet surtout de proposer des exercices.

\end{document}
